%% kapa: exam=1 round=36 year=2025 period=2 form=A subject=appraisal_law subject_ko=감정평가관계법규 seq=33 number=33
%% source: Q-Net 감정평가사 기출문제 (cst00309) artlSeq=5231525 file=raw/1008/5231525/2025년 제36회 감정평가사 1차시험 2교시 문제지 원본.pdf
%% conversion: kordoc 4.24.51 -> tools/gen1.py (auto); review=proofread-sonnet; flags=none
\begin{question}{33}
부동산등기법령상 등기의 신청에 관한 설명으로 옳지 않은 것은?
\begin{choices}
  \item 소유권보존등기 또는 소유권보존등기의 말소등기는 등기명의인으로 될 자 또는 등기명의인이 단독으로 신청한다.
  \item 공유물을 분할하는 판결에 의한 등기는 등기권리자 또는 등기의무자가 단독으로 신청한다.
  \item 가등기권리자는 가등기의무자의 승낙이 있을 때에는 단독으로 가등기를 신청할 수 있다.
  \item 채권자는 ｢민법｣에 따라 채무자를 대위하여 등기를 신청할 수 있다.
  \item 대표자나 관리인이 있는 법인 아닌 사단에 속하는 부동산의 등기에 관하여는 그 사단의 대표자나 관리인을 등기권리자 또는 등기의무자로 한다.
\end{choices}
\end{question}
\answer{5}
